
\section{Weighted angular-flux feasibility before source-matched stress closure}
\label{sec:weighted-flux-XXIV}
This section supplies a source-facing \emph{necessary} test for the
nonaxisymmetric correction route of Section~\ref{sec:nonaxis-Fourier-XXIII}.
The argument is elementary positivity and integration by parts; we do
not claim novelty for the covariance inequality or Navier--Stokes
energy identity. The proposed physical source and full unforced
blow-up theorem remain uncertified.

\begin{theorem}[Weighted covariance and two-slice torque obstruction]
\label{thm:weighted-flux-XXIV}
Let $v$ be a real smooth field, compactly supported in an annulus
$0<r<b$, periodic in $z$, and let
$Q_{ij}(r,z)=\langle v_i v_j\rangle_\theta$.
Define the azimuthal nonlinear mean momentum flux, its radial-slice
torque, and the two weighted diagonal moments by
\[
g_\theta=\frac1{r^2}\partial_r(r^2Q_{r\theta})
               +\partial_zQ_{z\theta},\quad
m(z)=\int_0^\infty r^2g_\theta(r,z)\,dr,
\]
\[
T(z)=\int_0^\infty r^2Q_{\theta z}(r,z)\,dr,\qquad
E_j(z)=\int_0^\infty r^2Q_{jj}(r,z)\,dr
\quad(j=\theta,z).
\]
Then, for every $z$,
\begin{equation}\label{eq:weighted-CS-XXIV}
m(z)=T'(z),\qquad
T(z)^2\le E_\theta(z)E_z(z),\qquad
2|T(z)|\le E_\theta(z)+E_z(z).
\end{equation}
For any primitive $M'(z)=m(z)$ and any interval $I$ in $z$,
\begin{equation}\label{eq:torque-osc-XXIV}
\boxed{\ \sup_{z\in I}(E_\theta(z)+E_z(z))
       \ge\operatorname{osc}_{z\in I}M(z).\ }
\end{equation}
If the radial support lies in $r\le b$ and
$\mathcal E_{\theta z}(z)=2\pi\int_0^\infty
r\,\langle v_\theta^2+v_z^2\rangle_\theta\,dr$ is the
non-halved directional kinetic-energy per axial slice, then
\begin{equation}\label{eq:physical-energy-XXIV}
\boxed{\ \sup_{z\in I}\mathcal E_{\theta z}(z)
 \ge \frac{2\pi}{b}\operatorname{osc}_{z\in I}M(z).\ }
\end{equation}
The bound is necessary for the \emph{portion of the actual source
assigned to the quadratic Reynolds stress}. In a full source
comparison all other linear, mean and pressure terms must first be
separated; $m$ cannot be identified with the entire source residual
without that decomposition.
\end{theorem}
\begin{proof}
Smoothness and radial compactness allow integration by parts in
$r$ and differentiation in $z$, giving $m=T'$.
The angular covariance matrix is positive semidefinite:
for any real $a,b$, $\langle(av_\theta+bv_z)^2\rangle_\theta\ge0$.
The resulting pointwise Cauchy--Schwarz bound, followed by weighted
integral Cauchy--Schwarz, gives $T^2\le E_\theta E_z$.
By arithmetic--geometric mean,
$2|T|\le E_\theta+E_z$.
For any two axial locations $z_+,z_-$,
\[
|M(z_+)-M(z_-)|=|T(z_+)-T(z_-)|
 \le |T(z_+)|+|T(z_-)|
 \le \sup_{z\in I}(E_\theta+E_z).
\]
Take the supremum over pairs, proving \eqref{eq:torque-osc-XXIV}.
Finally $E_\theta+E_z\le
b\mathcal E_{\theta z}/(2\pi)$ on $r\le b$, proving
\eqref{eq:physical-energy-XXIV}.
\end{proof}

\begin{corollary}[Exact Gaussian fixture and two-shell interpretation]
\label{cor:Gaussian-XXIV}
For the smooth $m=\pm1$ Gaussian model of
Proposition~\ref{prop:nonaxis-torque-XXIII},
\[
T(z)=\frac{\lambda}{8}\sin(2z),\qquad
m(z)=\frac{\lambda}{4}\cos(2z),\qquad
\operatorname{osc}_{[0,2\pi]}T=\frac{|\lambda|}{4}.
\]
Thus any stress realizing this \emph{entire} slice-torque target
must have $\sup_z(E_\theta+E_z)\ge|\lambda|/4$.
Conversely, the two-shell compensation of
Proposition~\ref{prop:two-shell-XXIII} makes $T\equiv0$ and
$m\equiv0$; it cannot reproduce a target with nonzero radial-slice
torque, even though its local angular source need not vanish.
\end{corollary}

\begin{proposition}[Zero-initial-energy exclusion for unforced smooth flow]
\label{prop:zero-data-XXIV}
If $u$ is a smooth finite-energy unforced solution on $\mathbb R^3$
or a periodic torus, and $u(0)=0$, then $u\equiv0$ on its entire
smooth lifespan. Therefore removing the force from the published
OpenAI construction \emph{while retaining its zero initial datum}
cannot prove a nontrivial unforced singularity.
\end{proposition}
\begin{proof}
The standard smooth energy identity yields
\[
\frac12\|u(t)\|_2^2+
\nu\int_0^t\|\nabla u(s)\|_2^2\,ds
=\frac12\|u(0)\|_2^2=0.
\]
Both left terms are nonnegative, so $u(t)=0$.
The originating forced theorem explicitly has $u(0)=0$
\cite{openai2026}. An unforced singularity attempt must instead
exhibit admissible \emph{nonzero} smooth initial data.
\end{proof}

\paragraph{Proof and release status.}
Independent Gaussian integration, $30{,}000$ rational positive
covariance tests and $30{,}000$ two-slice budget tests pass; a
separate scalar Lean module states the finite covariance inequality,
two-slice bound, and zero-data energy arithmetic. It does not
formalize angular averaging, Bochner integration, energy equality,
the source or the full PDE. The existing source certificate remains
\texttt{FAIL\_CLOSED\_MISSING\_SOURCE\_DATA}: source A.11/A.19
profiles, actual mean forcing, full cone, pressure, and zero-force
PDE residual are not authenticated. The seven-ray H8 realizability
and its spatially varying dual obstructions remain separate,
stronger gates. \emph{No unforced blow-up or complete Lean proof
is established.}
